% The Slide Definitions
%document
\documentclass[10pt]{beamer}
%theme
\usetheme{metropolis}
% packages
\usepackage{color}
\usepackage{listings}
\usepackage[ngerman]{babel}
\usepackage[utf8]{inputenc}
\usepackage{multicol}


% color definitions
\definecolor{mygreen}{rgb}{0,0.6,0}
\definecolor{mygray}{rgb}{0.5,0.5,0.5}
\definecolor{mymauve}{rgb}{0.58,0,0.82}

\lstset{
    backgroundcolor=\color{white},
    % choose the background color;
    % you must add \usepackage{color} or \usepackage{xcolor}
    basicstyle=\footnotesize\ttfamily,
    % the size of the fonts that are used for the code
    breakatwhitespace=false,
    % sets if automatic breaks should only happen at whitespace
    breaklines=true,                 % sets automatic line breaking
    captionpos=b,                    % sets the caption-position to bottom
    commentstyle=\color{mygreen},    % comment style
    % deletekeywords={...},
    % if you want to delete keywords from the given language
    extendedchars=true,
    % lets you use non-ASCII characters;
    % for 8-bits encodings only, does not work with UTF-8
    frame=single,                    % adds a frame around the code
    keepspaces=true,
    % keeps spaces in text,
    % useful for keeping indentation of code
    % (possibly needs columns=flexible)
    keywordstyle=\color{blue},       % keyword style
    % morekeywords={*,...},
    % if you want to add more keywords to the set
    numbers=left,
    % where to put the line-numbers; possible values are (none, left, right)
    numbersep=5pt,
    % how far the line-numbers are from the code
    numberstyle=\tiny\color{mygray},
    % the style that is used for the line-numbers
    rulecolor=\color{black},
    % if not set, the frame-color may be changed on line-breaks
    % within not-black text (e.g. comments (green here))
    stepnumber=1,
    % the step between two line-numbers.
    % If it's 1, each line will be numbered
    stringstyle=\color{mymauve},     % string literal style
    tabsize=4,                       % sets default tabsize to 4 spaces
    % show the filename of files included with \lstinputlisting;
    % also try caption instead of title
    language = [Sharp]C,
	showspaces = false,
	showtabs = false,
	showstringspaces = false,
	escapechar = ,
}

\def\ContinueLineNumber{\lstset{firstnumber=last}}
\def\StartLineAt#1{\lstset{firstnumber=#1}}
\let\numberLineAt\StartLineAt



\newcommand{\codeline}[1]{
	\alert{\texttt{#1}}
}


% Author and Course information
% This Document contains the information about this course.

% Authors of the slides
\author{\href{https://github.com/evilham}{@evilham}\\
\footnotesize{Basierend auf Folien von \href{https://github.com/satkowski}{@satkowski} (November 2016)}}

% Name of the Course
\institute{C\texttt{\#} Kurs}

% Fancy Logo
\titlegraphic{\hfill\includegraphics[height=1.25cm]{../templates/fsr_logo_cropped}}


% Presentation title
\title{Datenstrukturen}
\date{\today}


\begin{document}

\maketitle

\begin{frame}{Gliederung}
	\setbeamertemplate{section in toc}[sections numbered]
	\tableofcontents
\end{frame}

% ----------------------- Enummeration ------------------------------
\section{Enumeration}
\begin{frame}{Enumeration}
	\begin{itemize}
		\item Zu dt. Aufzählung
		\item Das Schlüsselwort dafür ist \alert{enum}
		\item Ist ein eigen erstellter Datentyp mit fester Anzahl konstanter Werte
		\item Den Werten werden Intern ein \alert{int}-Wert zugeordnet
		\begin{itemize}
			\item Diese beginnt Standardmäßig bei 0 und wird mit jedem neuen Element um 1 erhöht
			\item Diese Ordnung kann man auch selbst mit einem \alert{=} festlegen
			\item Man kann damit ein \alert{enum}-Wert in einen \alert{int}-Wert casten			
		\end{itemize}
		\item Syntax zur Nutzung:
		\begin{itemize}
			\item \texttt{\alert{<Enum> <Variablenname>} = \alert{<Enum>}.\alert{<Enumwert>};} oder
			\item \texttt{\alert{<Enum> <Variablenname>} = \alert{<Int-Wert>};}
		\end{itemize}
	\end{itemize}
\end{frame}

\begin{frame}{Enumeration}
	\lstinputlisting{resources/06_datenstrukturen/enum_creation.cs}
	\lstinputlisting{resources/06_datenstrukturen/enum_using.cs}
\end{frame}

% ----------------------- Structs ------------------------------
\section{Structs}
\begin{frame}{Structs}
	\begin{itemize}
		\item Schlüsselwort ist \alert{struct}
		\item Lässt sich genauso schreiben wie Klassen
		\item Sie sind nicht Referenz-, sondern Werttypen
		\item Werden genutzt um kleine Gruppen verwandter Variablen zusammen zufassen
		\item Sie können von Interfaces, aber nicht von anderen Structs oder Klassen erben 
	\end{itemize}
\end{frame}

\begin{frame}{Structs}
	\lstinputlisting{resources/06_datenstrukturen/struct_creation.cs}
	\lstinputlisting{resources/06_datenstrukturen/struct_using.cs}
\end{frame}

% ----------------------- Generics ------------------------------
\section{Generics}
\begin{frame}{Generics}
	\begin{itemize}
		\item Auch Typparameter genannt
		\item Ermöglicht es die Angabe von Typen in einer Klasse oder Methode zu verzögern
		\item Syntax:
		\begin{itemize}
			\item \texttt{class \alert{<Klassenname>} <\alert{<Typparameter>}> \{ \}}
		\end{itemize}
	\end{itemize}
	\textbf{Einschränkungen}\\
	\begin{itemize}
		\item Man kann den Typparameter auch einschränken
		\item Syntax:
		\begin{itemize}
			\item \texttt{\alert{<Klassenkopf>} where \alert{<Typparameter>} : \alert{<Einschränbkung>}\\ \{ \}}
		\end{itemize}
		\item Einschränkungen können sein:
		\begin{itemize}
			\item Basisklasse, Interfaces
			\item \alert{struct} (Werttyp), \alert{class} (Referenztyp)
			\item \alert{new()} (Typparameter muss einen öffentlichen, leeren Konstruktor haben)
		\end{itemize}
	\end{itemize}
\end{frame}

\begin{frame}{Generics}
	\lstinputlisting{resources/06_datenstrukturen/generics_creation.cs}
	\lstinputlisting{resources/06_datenstrukturen/generics_using.cs}
\end{frame}

% ----------------------- Collections ------------------------------
\section{Collections}
\begin{frame}{Collections}
	\begin{itemize}
		\item Sind in dem Framework vorhandene Klasse
		\item Dienen zur Auflistung von beliebigen Elementen
		\item Bieten Methoden zum Sortieren, Durchsuchen und Bearbeiten
		\begin{itemize}
			\item Nutzen dafür Generics
		\end{itemize}
		\item Haben eine dynamische Größe, können somit während der Ausführung schrumpfen und wachsen
		\item Namespace:
		\begin{itemize}			
			\item \alert{System.Collection.Generic}
		\end{itemize}
	\end{itemize}
\end{frame}

\subsection{List}
\begin{frame}{List}
	\begin{itemize}
		\item Eine Liste beliebiger Elemente
		\item Kann über einen Index angesteuert werden (ähnlich zum Array)
	\end{itemize}
	\lstinputlisting{resources/06_datenstrukturen/list.cs}
\end{frame}

\subsection{Dictionary}
\begin{frame}{Dictionary}
	\begin{itemize}
		\item Ein "Lexikon" von Paaren mit beliebigen Schlüssel und Wert
		\item Jeder Schlüssel kann nur genau einmal vorkommen
		\item Kann ebenfalls, ähnlich zum Array, angesprochen werden, aber dieses mal wird ein Schlüssel übergeben
	\end{itemize}
	\lstinputlisting{resources/06_datenstrukturen/dictionary.cs}
\end{frame}

\subsection{Weitere}
\begin{frame}{Weitere}
	\textbf{Stack}\\
	\begin{itemize}
		\item Ein ''Stapel'' an Elementen
		\item Arbeitet nach dem LIFO-Prinzip
		\item \alert{.Push()} (Hinzufügen eines Elementes), \alert{.Pop()} (Entfernen eines Elementes)
	\end{itemize}
	\lstinputlisting{resources/06_datenstrukturen/stack.cs}
	\textbf{Queue}\\
	\begin{itemize}
		\item Eine ''Schlange'' an Elementen
		\item Arbeitet nach dem FIFO-Prinzip
		\item \alert{.Enqueue()} (Hinzufügen eines Elementes), \alert{.Dequeue()} (Entfernen eines Elementes)
	\end{itemize}
	\lstinputlisting{resources/06_datenstrukturen/queue.cs}
\end{frame}

% ----------------------- Indexer ------------------------------
\section{Indexer}
\begin{frame}{Indexer}
	\begin{itemize}
		\item Sind Member einer Klasse
		\item Geben die Möglichkeit wie ein/e Array/Collection über \alert{[ ]} angesprochen zu werden
		\item Syntax:
		\begin{itemize}
			\item \texttt{\alert{<Typ>} this[\alert{<Indextyp> <Indexname>}]\\ \{ get \{ \alert{<Code>} \} set \{ \alert{<Code>}\} \}}
		\end{itemize}
	\end{itemize}
\end{frame}

\begin{frame}{Indexer}
	\lstinputlisting{resources/06_datenstrukturen/indexer.cs}
\end{frame}

\end{document}
